\documentclass[11pt]{article}

\usepackage[T1]{fontenc}
\usepackage[utf8]{inputenc}
\usepackage{lmodern}
\usepackage{microtype}
\usepackage[margin=1in]{geometry}
\usepackage{amsmath,amssymb,mathtools}
\usepackage{booktabs,tabularx,array}
\usepackage{xcolor}
\usepackage{enumitem}
\usepackage{fancyhdr}
\usepackage{hyperref}

\definecolor{linkblue}{HTML}{245A8D}
\definecolor{rulegray}{HTML}{8A9299}
\hypersetup{
  colorlinks=true,
  linkcolor=linkblue,
  citecolor=linkblue,
  urlcolor=linkblue,
  pdftitle={Smooth Collapse of Trapped Void Fragments in an Incompressible VOF Method},
  pdfsubject={Revision 3: capillary-biased signed compliance}
}

\setlist{nosep,leftmargin=*}
\setlength{\parindent}{0pt}
\setlength{\parskip}{0.55em}
\setlength{\headheight}{14pt}
\pagestyle{fancy}
\fancyhf{}
\fancyhead[L]{Trapped-void collapse in VOF}
\fancyhead[R]{Working research note}
\fancyfoot[C]{\thepage}
\renewcommand{\headrulewidth}{0.4pt}

\newcommand{\dd}{\mathop{}\!\mathrm{d}}
\newcommand{\Dt}{\Delta t}
\newcommand{\pos}[1]{\left[#1\right]_+}
\newcommand{\divg}{\nabla\!\cdot}
\newcommand{\uv}{\boldsymbol{u}}
\newcommand{\trapped}{\mathcal{T}}
\newcommand{\doi}[1]{\href{https://doi.org/#1}{doi:#1}}

\title{\vspace{-2.2em}\textbf{Smooth Collapse of Trapped Void Fragments\\
in an Incompressible Volume-of-Fluid Method}\\[0.6em]
\large A working research note motivated by Truchas}
\author{}
\date{15 September 2026}

\begin{document}
\maketitle
\vspace{-2.4em}

\begin{abstract}
In a one-fluid incompressible volume-of-fluid (VOF) calculation, a resolved void region
provides a prescribed void-pressure boundary to the liquid.  When a small void fragment
becomes enclosed within one or a few mixed cells, that internal boundary is no longer
represented.  The mixed cells are then treated as incompressible and a nonzero residual
void fraction may persist indefinitely.  Truchas contains an experimental collapse model
whose rate is tied to the magnitude of the recent local change in void fraction.  This
adapts collapse to the local interface-evolution time scale, but it switches itself off
when the VOF field becomes stationary.  Pressure-driven filling against a closed ceiling
is a particularly clear failure case.

This note organizes several possible remedies.  They include fixed-time relaxation,
flow-time scaling, pressure-driven relaxation, Rayleigh-collapse-inspired laws, a
compressible trapped-gas model, and a void-fraction-dependent pressure compliance.  The
last of these gives a useful synthesis: an enclosed mixed cell is treated as a
continuously disappearing internal pressure boundary.  It adds a nonnegative reaction
term to the pressure projection, anchors a nearly void cell near the void pressure, and
recovers ordinary incompressibility as the void fraction vanishes.  The note derives the
common kinematic framework, gives both absolute-pressure and incremental
pressure-correction forms, and develops a capillary-biased signed variant.  That variant
allows a large unresolved pocket to expand or contract like open void, but increasingly
biases a small remnant toward collapse through a subgrid Young--Laplace pressure shift.
The note identifies unresolved choices and proposes a staged implementation and
verification program.
\end{abstract}

\hrule height 0.6pt
\vspace{0.4em}
\textbf{Status.} Revision 3 adds capillary-biased signed compliance and its Truchas
incremental pressure-correction form in Sec.~6.  A first signed-compliance prototype has
performed well in limited testing; the capillary bias remains a candidate extension.
This is a research note, not a specification.  Equations labeled as
candidate models are hypotheses to be evaluated.  The reconstruction of the existing
Truchas algorithm should be checked against the precise source revision used for any
implementation work.
\vspace{0.4em}
\hrule height 0.6pt

\clearpage
\tableofcontents
\clearpage

\section{Problem and modeling scope}

Let \(F\) denote the volume fraction of real incompressible fluid in a VOF
description \cite{hirt-nichols1981}, and let
\(\alpha\) denote void fraction.  For now assume that no solid fraction occupies the
cell, so
\begin{equation}
  F+\alpha=1.
\end{equation}
A completely void cell adjacent to liquid acts as part of the free surface and supplies
a prescribed pressure, denoted \(p_v\) (often zero gauge pressure), to the liquid pressure
solve.  Once void is enclosed in mixed cells and no cell in the fragment is treated as
pure void, the discrete internal \(p=p_v\) boundary disappears.  Ordinary projection
then makes the mixed region incompressible.  The remaining void behaves numerically like
an incompressible constituent and can persist even though the application intends void
to be an absence of material rather than a conserved gas phase.

The immediate objective is therefore not a complete physical model of gas bubbles.  It
is a subgrid model that permits an unresolved, enclosed void volume to be replaced by
liquid through actual liquid flux.  The model should not act on the exterior-connected
free surface and should turn itself off continuously when the void has vanished.

Two distinct goals must remain separate:
\begin{enumerate}
  \item A \emph{numerical void-annihilation model} intentionally permits enclosed void
        volume to reach zero.
  \item A \emph{compressible trapped-gas model} conserves gas mass.  Compression changes
        its volume but, without dissolution, condensation, or venting, does not make it
        disappear.
\end{enumerate}
The first goal is the one addressed primarily here.

\section{Kinematic framework: collapse as permitted volumetric strain}

Conservation of incompressible real fluid is
\begin{equation}
  \frac{\partial F}{\partial t}+\divg(F\uv)=0.
  \label{eq:fluid-conservation}
\end{equation}
Expanding the divergence and using \(F=1-\alpha\) gives
\begin{equation}
  \frac{D\alpha}{Dt}=F\,\divg\uv=(1-\alpha)\divg\uv.
  \label{eq:kinematic-identity}
\end{equation}
Suppose an eligible trapped fragment is assigned a nonnegative void-destruction rate
\(R\):
\begin{equation}
  \frac{D\alpha}{Dt}=-R(\alpha,p,\uv,\ldots), \qquad R\geq 0.
\end{equation}
Consistency with Eq.~\eqref{eq:fluid-conservation} requires the modified divergence
constraint
\begin{equation}
  \boxed{\divg\uv=-\frac{R}{1-\alpha}.}
  \label{eq:general-divergence}
\end{equation}
Thus trapped-void collapse is naturally formulated as a localized negative volumetric
strain.  Liquid replacing the lost void must enter through the velocity fluxes; the
volume fraction is not simply edited after the flow solve.  Equation
\eqref{eq:general-divergence} is the common framework for all rate laws considered below.

The factor \((1-\alpha)^{-1}\) is kinematically required but becomes singular in the
pure-void limit.  That endpoint is already governed by the pure-void pressure boundary,
so Eq.~\eqref{eq:general-divergence} is intended for mixed cells.  A discrete method must
join these regimes carefully.

\section{The existing Truchas model}

The Truchas manual describes an experimental model intended to eliminate persistent
void fragments that VOF otherwise treats as incompressible.  It uses a dimensionless
relaxation factor in \([0,1]\), with default \(0.1\), and describes the factor as roughly
the inverse of the number of time steps over which inertial forces are allowed to
collapse a cell's void \cite{truchas-manual}.  The model is implemented in the pressure
projection, which is the appropriate algorithmic location for enforcing a nonzero
divergence \cite{bell-colella-glaz1989}.

\subsection{Source-level reconstruction}

The source-level behavior discussed here can be summarized as follows.  Define the
ordinary pre-collapse VOF change
\begin{equation}
  \Delta\alpha=\alpha^{n+1,*}-\alpha^n.
\end{equation}
For an enclosed mixed cell, the collapse amount made available to the projection is,
schematically,
\begin{equation}
  \boxed{\delta\alpha_c=\min\!\left(|\Delta\alpha|,\,r\alpha\right),}
  \label{eq:existing-rule}
\end{equation}
where \(r\) is \texttt{void\_collapse\_relaxation}.  Both signs of
\(\Delta\alpha\) contribute the same collapse sign; consequently the method does not
merely continue an already shrinking fragment.  It uses \(|\Delta\alpha|\) as a sensor
for the rate at which the local VOF configuration is evolving.  The exact algebra and
candidate-cell test should be cited to a recorded Truchas commit before this note is
circulated as a code description \cite{truchas-source}.

\subsection{Plausible motivation}

When the cap in Eq.~\eqref{eq:existing-rule} is inactive,
\begin{equation}
  \frac{\delta\alpha_c}{\Dt}\simeq
  \left|\frac{D\alpha}{Dt}\right|.
\end{equation}
The implied local VOF-evolution time is
\begin{equation}
  \tau_\alpha=\frac{\alpha}{|D\alpha/Dt|}.
\end{equation}
The collapse therefore proceeds on approximately the same intrinsic time scale as the
resolved interface evolution.  For an interface sweeping across cells,
\(|D\alpha/Dt|\sim U/h\), so the scale is related to a cell-crossing time.  This is a
clever heuristic: it avoids prescribing an independent dimensional collapse time,
responds rapidly in violent filling and splashing, and is approximately invariant to
time-step refinement while the uncapped estimate is valid.  The parameter \(r\) then
acts mainly as a per-step robustness limiter.

\subsection{Failure when the VOF field becomes stationary}

The same sensor creates the central defect:
\begin{equation}
  |\Delta\alpha|\longrightarrow0
  \quad\Longrightarrow\quad
  \delta\alpha_c\longrightarrow0,
\end{equation}
even when \(\alpha>0\) and the fragment remains enclosed.  The model is triggered by
recent kinematic evolution, not by the present existence of trapped void or by the force
driving its collapse.

A pressure-fed mold filling against a closed horizontal ceiling is an almost diagnostic
counterexample.  Once the advancing liquid reaches the ceiling, the interface stops
moving through the grid and the local VOF field approaches a stationary state.  The
remaining top-layer void can still be substantial, and the liquid pressure may remain
positive, but Eq.~\eqref{eq:existing-rule} switches off.  Ordinary incompressibility then
prevents further net inflow and freezes the residual void.

There is a second conceptual weakness.  A translating or deforming fragment can produce
positive \(\Delta\alpha\) in some cells and negative \(\Delta\alpha\) in others without
changing its total void volume.  Local \(|\Delta\alpha|\) is therefore not, in general,
a measure of component collapse.

\section{Requirements for a successor model}

A useful replacement should satisfy the following requirements.
\begin{enumerate}
  \item \textbf{Persistent activation:} an eligible cell with \(\alpha>0\) must not
        become inactive merely because \(\Delta\alpha=0\).
  \item \textbf{Flux consistency:} lost void is replaced by liquid through the velocity
        field, using Eq.~\eqref{eq:general-divergence}.
  \item \textbf{Free-surface protection:} exterior-connected void must remain an
        ordinary free surface rather than being annihilated.
  \item \textbf{Correct endpoints:} the pure-void limit behaves like \(p=p_v\), while
        the zero-void limit recovers \(\divg\uv=0\).
  \item \textbf{Pressure response:} positive pressure loading should continue to drive
        collapse in a stagnating pocket.
  \item \textbf{Time-step convergence:} parameters should define a continuous-time law;
        per-step bounds may remain as numerical limiters.
  \item \textbf{Controlled grid dependence:} any dependence on cell size or inferred
        fragment size should be explicit and testable.
  \item \textbf{Solver compatibility:} the pressure operator should preferably remain
        coercive, or any nonlinearity should admit a robust active-set or fixed-point
        treatment.
\end{enumerate}

\section{Candidate collapse-rate models}

Table~\ref{tab:candidates} previews the alternatives.  None of the constitutive laws by
itself settles the separate question of which void is eligible for collapse.

\begin{table}[htbp]
\centering
\small
\caption{Candidate laws for \(D\alpha/Dt=-R\).}
\label{tab:candidates}
\begin{tabularx}{\textwidth}{@{}>{\raggedright\arraybackslash}p{0.18\textwidth}
  >{\raggedright\arraybackslash}p{0.27\textwidth}X@{}}
\toprule
Model & Rate \(R\) & Principal consequence \\
\midrule
Existing Truchas &
  \(\min(|\Delta\alpha|,r\alpha)/\Dt\) &
  Flow-responsive and simple, but becomes identically inactive for stationary trapped
  void. \\
Fixed time &
  \(\alpha/\tau_c\) &
  Persistent exponential decay with a clear continuous-time parameter; no pressure or
  local-dynamics response. \\
Local flow time &
  \(C_u\alpha |\uv|/h\) or \(C_D\alpha\|\boldsymbol D\|\) &
  Uses resolved dynamics without a dimensional time, but either reacts to translation
  or fails in pressure-loaded stagnation. \\
Pressure relaxation &
  \(\alpha\,g((p-p_v)/p_0)/\tau_0\) &
  Continues under static pressure loading; linear or rate-saturated variants are easy to
  construct. \\
Rayleigh-inspired &
  \(K\alpha^{2/3}\sqrt{\pos{p-p_v}/\rho}/h\) in 3-D &
  Inertial, pressure- and size-responsive, and reaches zero in finite time; introduces
  geometric and grid/component-size assumptions. \\
Compressible gas &
  EOS plus gas-mass conservation &
  Physically interpretable pressure-volume response, but gas does not disappear without
  additional mass-transfer physics. \\
Capillary-biased compliance &
  \(\divg\uv=-C(\alpha)[p-p_v+\Pi_\sigma(\alpha,h)]\) &
  Retains reversible open-void behavior for large fragments while shifting small
  remnants toward collapse; remains linear when the capillary shift is lagged. \\
\bottomrule
\end{tabularx}
\end{table}

\subsection{Fixed-time relaxation}

The simplest persistent law is
\begin{equation}
  R=\frac{\alpha}{\tau_c}, \qquad
  \divg\uv=-\frac{\alpha}{(1-\alpha)\tau_c}.
  \label{eq:fixed-time}
\end{equation}
It gives exponential decay along a trajectory and remains active until the void becomes
negligible.  A per-step bound such as \(\delta\alpha_c\le r_{\max}\alpha\) can be retained
for robustness, but should not define the constitutive law.  This model is an excellent
baseline because its time-step convergence is easy to verify.  Its weakness is the new
dimensional parameter and its indifference to pressure: it collapses eligible void at the
same rate whether weakly or strongly loaded.

Choosing \(\tau_c=N_c\Dt\) reproduces a fraction-per-step behavior but forfeits a
continuous-time model.  It may be useful only as a temporary compatibility mode.

\subsection{Resolved-flow time scales}

The apparent intent of the existing method can be preserved by replacing the vanishing
VOF-change sensor with a nonzero local flow scale, for example
\begin{align}
  R_u &= C_u\alpha\frac{|\uv|}{h}, \\
  R_D &= C_D\alpha\|\boldsymbol D(\uv)\|,
  \qquad
  \boldsymbol D(\uv)=\tfrac12(\nabla\uv+\nabla\uv^T).
\end{align}
The velocity scale collapses a fragment that is merely translating, even if it is not
being compressed, and it vanishes at the stagnating ceiling.  The strain-rate scale
avoids sensitivity to rigid translation, but it too can vanish in a pressure-loaded
stagnant pocket.  These scales may be useful modifiers or fallback estimates, but neither
addresses the motivating failure by itself.

\subsection{Pressure-driven relaxation}

A simple pressure-responsive family is
\begin{equation}
  R=\frac{\alpha}{\tau_0}
  g\!\left(\frac{p-p_v}{p_0}\right),
  \label{eq:pressure-relaxation}
\end{equation}
where \(g(x)\ge0\) for \(x>0\).  Two useful choices are
\begin{align}
  g(x)&=\pos{x}, &&\text{linear pressure response},\\
  g(x)&=\min(\pos{x},1), &&\text{saturated response}.
\end{align}
The linear choice is equivalently \(R=M\alpha\pos{p-p_v}\), with
\(M=(\tau_0p_0)^{-1}\).  The saturated form defines \(\tau_0\) as the fastest relaxation
time and prevents the rate from becoming arbitrarily large at high pressure.  Both remain
active in a stationary, pressure-loaded ceiling pocket.  They introduce a pressure scale
and a collapse time, but those parameters have clear units and meanings.

The one-sided positive part suppresses void growth when \(p<p_v\).  Omitting it produces
a reversible compliance model that can grow void under tension, which would amount to a
rudimentary cavitation model and should be a separate decision.

\subsection{Rayleigh-collapse-inspired inertial scaling}

For an empty spherical cavity of initial radius \(R_0\) in an infinite inviscid liquid
under constant pressure difference \(\Delta p>0\), Rayleigh's collapse time is
\begin{equation}
  t_R\simeq0.915R_0\sqrt{\frac{\rho}{\Delta p}}
  \label{eq:rayleigh-time}
\end{equation}
\cite{rayleigh1917,brennen1995}.  This suggests the intrinsic scale
\begin{equation}
  \tau_p=C_R L\sqrt{\frac{\rho}{\pos{p-p_v}}}.
  \label{eq:pressure-time}
\end{equation}
Taking \(L=h\) yields a deliberately cell-scale model.  Taking \(L\) as the radius of an
entire connected component is more physically meaningful but requires component
identification and distributed reductions.

At the cell level in three dimensions, an effective spherical radius is
\begin{equation}
  R_{\mathrm{eff}}=
  \left(\frac{3\alpha V_c}{4\pi}\right)^{1/3}.
\end{equation}
Combining \(R\sim\alpha/\tau_p\) with this radius gives the scaling
\begin{equation}
  R_{\mathrm{inertial}}=
  K_3\frac{\sqrt{\pos{p-p_v}/\rho}}{V_c^{1/3}}\alpha^{2/3}.
  \label{eq:inertial-rate}
\end{equation}
In two dimensions the analogous exponent is \(1/2\).  Because
\(\dot\alpha\propto-\alpha^{2/3}\), the three-dimensional model reaches zero in finite
time.  The coefficient must be calibrated carefully: using the instantaneous Rayleigh
time as \(\alpha/R\) makes the total extinction time three times the initial local time,
so matching Eq.~\eqref{eq:rayleigh-time} requires the corresponding factor of three.

Equation~\eqref{eq:inertial-rate} is not the Rayleigh--Plesset equation.  It borrows its
pressure, density, length, and finite-extinction scaling for a subgrid closure.  Cell
geometry, walls, viscosity, surface tension, nonsphericity, and pressure variation are
all omitted.  It is therefore a second-generation candidate, not the first model to
implement.

\subsection{Actual compressible trapped gas}

If the void is physically entrapped air, a more faithful model conserves gas mass and
uses an equation of state, for example
\begin{equation}
  p_g V_g^\gamma=\text{constant}.
\end{equation}
Li and Carrica tightly couple bubble density, void fraction, pressure, and velocity in a
two-phase solver with incompressible liquid and compressible bubbles
\cite{li-carrica2018}.  Shin uses a polytropic gas law and represents gas compressibility
in the source term of a unified pressure Poisson equation \cite{shin2020}.  These works
support the general numerical idea of coupling pressure and void-volume response
implicitly.

They do not, however, solve the present annihilation problem: conserved gas can become
highly compressed but cannot vanish.  Dissolution, condensation, leakage, or venting
would be required.  A compressible-gas model should therefore be treated as a distinct
future physics capability, not as a more elaborate implementation of numerical cleanup.

\section{Pressure compliance as a disappearing internal boundary}

The most useful synthesis is to interpret an unresolved trapped fragment as an internal
\(p=p_v\) boundary that should disappear continuously as \(\alpha\to0\).  Introduce a
nonnegative pressure compliance \(C(\alpha,p)\) through
\begin{equation}
  \boxed{\divg\uv=-C(\alpha,p)\,(p-p_v).}
  \label{eq:compliance}
\end{equation}
For a one-sided annihilation law, replace \(p-p_v\) by \(\pos{p-p_v}\).

The desired limits are
\begin{equation}
  C\longrightarrow\infty \quad\text{as}\quad \alpha\to1^{-},
  \qquad
  C\longrightarrow0 \quad\text{as}\quad \alpha\to0^{+}.
  \label{eq:compliance-limits}
\end{equation}
A large compliance strongly anchors the pressure near \(p_v\) while permitting rapid
net inflow.  A vanishing compliance recovers ordinary incompressibility and lets the
pressure approach that of the surrounding liquid.  More precisely, the model provides a
large \emph{capacity} for negative divergence near the void limit; the actual divergence
is determined by the coupled global pressure solve.

The qualitative feedback loop is
\begin{align*}
  \text{enclosed void}
  &\;\Rightarrow\;
  \text{pressure compliance}
  \;\Rightarrow\;
  \text{negative divergence}\\
  &\;\Rightarrow\;
  \text{liquid inflow}
  \;\Rightarrow\;
  \text{less void and less compliance}.
\end{align*}
This loop does not require an antecedent \(\Delta\alpha\).

\subsection{Connection to a void-fraction rate law}

If
\begin{equation}
  \frac{D\alpha}{Dt}=-\lambda(\alpha)\pos{p-p_v},
\end{equation}
then Eq.~\eqref{eq:kinematic-identity} gives
\begin{equation}
  C(\alpha)=\frac{\lambda(\alpha)}{1-\alpha}.
\end{equation}
The linear pressure-relaxation model is obtained with
\begin{equation}
  \lambda(\alpha)=\frac{\alpha}{\tau_0p_0},
  \qquad
  \boxed{C(\alpha)=
  \frac{\alpha}{(1-\alpha)\tau_0p_0}.}
  \label{eq:simple-compliance}
\end{equation}
This choice simultaneously gives exponential void decay at fixed pressure, strong
anchoring near the void endpoint, and exact recovery of incompressibility as
\(\alpha\to0\).

More generally,
\begin{equation}
  C(\alpha)=C_0\frac{\alpha^m}{(1-\alpha)^n},
  \qquad m>0,\;n>0,
\end{equation}
allows the two endpoint transitions to be shaped independently.  There is little reason
to introduce both exponents before the simplest choice \(m=n=1\) has been tested.

\subsection{Signed, one-sided, and capillary-biased response}

The symmetric form of Eq.~\eqref{eq:compliance},
\begin{equation}
  \divg\uv=-C(\alpha)(p-p_v),
  \label{eq:signed-compliance}
\end{equation}
allows both signs of divergence.  It collapses an existing fragment when \(p>p_v\) and
expands it when \(p<p_v\).  By contrast,
\begin{equation}
  \divg\uv=-C(\alpha)\pos{p-p_v}
  \label{eq:one-sided-compliance}
\end{equation}
is an irreversible annihilation law.  The signed law is linear in pressure and a first
prototype has worked well in limited tests.  Moreover, positive divergence in a cell
that is mostly void need not be interpreted as nucleating a new cavity: it can represent
motion of an already existing unresolved liquid--void interface.  If \(C(0)=0\), a
completely liquid cell cannot nucleate void through this closure alone.

A compromise is to retain signed compliance but shift its equilibrium pressure by a
capillary term:
\begin{equation}
  \boxed{
  \divg\uv=-C(\alpha)
  \left[p-p_v+\Pi_\sigma(\alpha,h)\right]
  =-C(\alpha)\left[p-p_{\mathrm{eq}}\right],}
  \qquad
  p_{\mathrm{eq}}=p_v-\Pi_\sigma.
  \label{eq:capillary-signed}
\end{equation}
Thus the fragment collapses for \(p>p_v-\Pi_\sigma\), is in equilibrium at
\(p=p_v-\Pi_\sigma\), and expands only when the liquid pressure falls below that
threshold.  For a large, nearly open void, \(\Pi_\sigma\) should tend to zero so that
both signs remain available around \(p_v\).  As the unresolved fragment becomes small,
\(\Pi_\sigma\) increases and progressively suppresses expansion.

The sign follows the Young--Laplace jump.  For a spherical void in three dimensions,
\begin{equation}
  p_v-p_{\ell}=\frac{2\sigma}{R_{\mathrm{eff}}},
  \label{eq:young-laplace}
\end{equation}
so the liquid-side equilibrium pressure is below the void pressure.  More generally the
curvature scale is \((d-1)/R_{\mathrm{eff}}\) in \(d\) spatial dimensions
\cite{brackbill1992}.  A deliberately cell-scale closure is
\begin{equation}
  \boxed{
  \Pi_\sigma(\alpha,h)
  =c_\sigma\frac{\sigma}{h}B(\alpha),}
  \qquad
  h=\begin{cases}
       V_c^{1/3}, & \text{three dimensions},\\
       A_c^{1/2}, & \text{two dimensions},
     \end{cases}
  \label{eq:capillary-shift}
\end{equation}
where \(\sigma\) is the physical liquid--void surface tension, \(c_\sigma\) is a
dimensionless subgrid curvature/geometry coefficient, and \(B(\alpha)\) controls when
the bias turns on.  A fixed, documented \(B\) gives this part of the model two inputs,
\(\sigma\) and \(c_\sigma\).  If \(\sigma\) is already a material property, only
\(c_\sigma\) is a new calibration parameter.  If it is not independently known, the
projection identifies only their product
\(\Sigma_{\mathrm{eff}}=c_\sigma\sigma\).  Any adjustable transition point or exponent
inside \(B\) would add parameters and should be avoided initially.

A parameter-free baseline shape that satisfies the required open-void limit is
\begin{equation}
  \boxed{B(\alpha)=(1-\alpha)^2.}
  \label{eq:baseline-bias-shape}
\end{equation}
It gives \(B(0)=1\), \(B(1)=0\), and, together with
Eq.~\eqref{eq:simple-compliance}, \(C B\sim\alpha(1-\alpha)\).  A more literal
effective-radius model could multiply this fade by \(\alpha^{-1/d}\), but that introduces
a singular endpoint and a regularization choice.  The bounded quadratic shape is the
cleaner first test.

The limiting behavior of the \emph{product} \(C\Pi_\sigma\), not just
\(\Pi_\sigma\), must be checked.  Since Eq.~\eqref{eq:simple-compliance} has
\(C\sim(1-\alpha)^{-1}\) as \(\alpha\to1\), recovering an unbiased open-void limit
requires \(B(\alpha)=o(1-\alpha)\) there, or an explicit cap/handoff must keep
\(C\Pi_\sigma\) bounded and vanishing.  The closure should be applied only to isolated,
underresolved fragments.  If resolved surface tension already supplies the relevant
pressure jump, the subgrid term must be blended or disabled to avoid double counting.

\subsection{Absolute-pressure and incremental pressure-correction forms}

The compliance model is compatible with either a projection that solves for the new
pressure or an incremental method that uses an existing pressure in the predictor and
solves for a correction.  The distinction changes the algebraic right-hand side, not the
underlying model.

\subsubsection{Absolute-pressure form}

Let \(\widetilde{\uv}\) be a provisional velocity that has not yet received the new
pressure force, and write
\begin{equation}
  \uv^{n+1}=\widetilde{\uv}-\boldsymbol{K}\nabla p^{n+1},
  \qquad
  \boldsymbol{K}=\Dt\,\rho^{-1}\boldsymbol{I}
  \label{eq:absolute-update}
\end{equation}
for the simplest scalar-density projection.  Here \(\boldsymbol K\) may instead denote
the actual discrete face mobility used by the code.  With compliance lagged in
\(\alpha\), substitution into Eq.~\eqref{eq:compliance} gives
\begin{equation}
  \boxed{
  -\divg\!\left(\boldsymbol{K}\nabla p^{n+1}\right)
  +C(\alpha^*)p^{n+1}
  =-\divg\widetilde{\uv}+C(\alpha^*)p_v.}
  \label{eq:reaction-poisson}
\end{equation}
Thus an eligible mixed cell receives a nonnegative diagonal reaction term.
Equation~\eqref{eq:reaction-poisson} is a smooth penalty or Robin-like transition:
large \(C\) approximates a void-pressure condition, intermediate \(C\) partially anchors
pressure, and \(C=0\) gives the ordinary interior projection.  The added reaction term is
coercive and should improve rather than damage the definiteness of the linear operator,
although strong coefficient variation can still affect conditioning.

For capillary-biased signed compliance, replace \(p_v\) by
\(p_{\mathrm{eq}}=p_v-\Pi_\sigma\).  With \(\alpha\), \(h\), and \(\Pi_\sigma\) lagged,
the absolute-pressure equation remains linear:
\begin{equation}
  -\divg(\boldsymbol K\nabla p^{n+1})+Cp^{n+1}
  =-\divg\widetilde{\uv}+C(p_v-\Pi_\sigma).
  \label{eq:capillary-absolute}
\end{equation}

\subsubsection{Incremental pressure-correction form}

For Truchas, let \(p^n\) be the existing pressure used in the predictor.  The incremental
method can be written
\begin{align}
  \uv^* &= \widetilde{\uv}-\boldsymbol{K}\nabla p^n,
  \label{eq:incremental-predictor}\\
  \uv^{n+1} &= \uv^*-\boldsymbol{K}\nabla\phi,
  \label{eq:incremental-velocity}\\
  p^{n+1} &= p^n+\phi,
  \label{eq:incremental-pressure}
\end{align}
where \(\phi\) is the pressure correction.  Requiring the corrected velocity to satisfy
\begin{equation}
  \divg\uv^{n+1}=-C(\alpha^*)\bigl(p^{n+1}-p_v\bigr)
\end{equation}
gives the incremental analog of Eq.~\eqref{eq:reaction-poisson}:
\begin{equation}
  \boxed{
  -\divg\!\left(\boldsymbol{K}\nabla\phi\right)
  +C(\alpha^*)\phi
  =-\divg\uv^*-C(\alpha^*)\bigl(p^n-p_v\bigr).}
  \label{eq:reaction-poisson-incremental}
\end{equation}
This is the form directly applicable to Truchas.  The reaction coefficient multiplies
the pressure \emph{correction}, but the right-hand side must also contain the offset of
the predictor pressure from the void pressure.  Omitting that offset would impose
\(\divg\uv^{n+1}=-C\phi\), anchoring the correction to zero rather than anchoring the
updated physical pressure to \(p_v\).

Equations~\eqref{eq:reaction-poisson} and
\eqref{eq:reaction-poisson-incremental} are algebraically equivalent when the predictor
and corrector use the same discrete gradient, divergence, mobility, and pressure boundary
treatment.  Substituting \(\phi=p^{n+1}-p^n\) and
Eq.~\eqref{eq:incremental-predictor} into
Eq.~\eqref{eq:reaction-poisson-incremental} recovers
Eq.~\eqref{eq:reaction-poisson}.  Incremental correction therefore does not obstruct the
proposed compliance model; it simply splits the absolute-pressure equation around the
pressure already used in the predictor.

The corresponding capillary-biased incremental equation is
\begin{equation}
  \boxed{
  -\divg(\boldsymbol K\nabla\phi)+C\phi
  =-\divg\uv^*-C\left[p^n-p_v+\Pi_\sigma(\alpha^*,h)\right].}
  \label{eq:capillary-incremental}
\end{equation}
Only the right-hand-side offset differs from
Eq.~\eqref{eq:reaction-poisson-incremental}; the matrix and its nonnegative reaction
term are unchanged.  Consequently the capillary-biased signed model preserves the
linear pressure-correction solve when \(\Pi_\sigma\) is evaluated from lagged geometry.

If the code uses a scaled pressure update, for example
\(p^{n+1}=p^n+\omega_p\phi\), the reaction contribution multiplying \(\phi\) becomes
\(\omega_p C\phi\).  Likewise, \(\boldsymbol K\) must be the actual coefficient relating
the velocity correction to \(\nabla\phi\), including any BDF or variable-density factors.
These details should be taken from the Truchas discrete update rather than inferred from
the schematic equations above.

Using \(\pos{p-p_v}\) makes the strict one-sided equation semilinear.  A simple active-set method can
classify cells from the trial updated pressure \(p^n+\phi\), solve the resulting linear
system, and repeat until the active set stops changing.  In incremental form, with
\(A_i=1\) on the active set and zero otherwise, Eq.~\eqref{eq:reaction-poisson-incremental}
becomes
\begin{equation}
  -\divg(\boldsymbol K\nabla\phi)+A C\phi
  =-\divg\uv^*-A C(p^n-p_v).
  \label{eq:active-incremental}
\end{equation}
A smooth approximation to the positive part is another option, but it introduces a
regularization scale and permits some response for \(p<p_v\).

\subsection{Endpoint treatment and discrete regularization}

The singular compliance in Eq.~\eqref{eq:simple-compliance} is a meaningful limiting
statement, not a number to evaluate at \(F=0\).  Practical options include:
\begin{enumerate}
  \item retain the existing pure-void Dirichlet treatment for \(F\le F_D\), and use
        compliance only for \(F>F_D\);
  \item cap \(C\) at a value that makes the reaction term dominate the local pressure
        stencil by a prescribed factor;
  \item replace \(1-\alpha\) by \(\max(1-\alpha,F_\epsilon)\), while verifying that
        results are insensitive to \(F_\epsilon\);
  \item integrate the local linear pressure-relaxation law exactly over a step, reducing
        sensitivity to a large explicit collapse rate.
\end{enumerate}
The first two preserve the interpretation most transparently.  A smooth handoff can be
constructed by choosing the cap to match the discrete Dirichlet penalty at \(F_D\).

\section{Eligibility: which void is allowed to collapse?}

The rate law and the eligibility test are logically independent.  The ideal criterion is
topological:
\begin{equation}
  \trapped_i=
  \begin{cases}
    1, & \text{void in cell \(i\) is disconnected from exterior void},\\
    0, & \text{otherwise}.
  \end{cases}
\end{equation}
Only cells with \(\trapped_i=1\) receive the modified divergence constraint.

The existing inexpensive criterion---a mixed cell with no neighboring completely void
cell---captures many one-cell fragments but can misclassify a thin or underresolved
free-surface filament.  It also cannot identify a trapped component spanning several
mixed cells when no member crosses the pure-void threshold.

A graph-based method would mark cells carrying appreciable void, connect face neighbors,
seed components touching exterior void or a pressure outlet, and classify all unseeded
components as trapped.  In parallel this becomes a distributed flood fill or
connected-component problem.  It need not necessarily be run every step: connectivity
could be updated only after material topology changes beyond a threshold.  Hysteresis in
the graph thresholds is desirable to prevent classification chatter.

Component identification also enables physically better length scales and diagnostics:
\begin{equation}
  V_{v,k}=\sum_{i\in k}\alpha_i V_i,
  \qquad
  R_{k}=\left(\frac{3V_{v,k}}{4\pi}\right)^{1/3}.
\end{equation}
However, global connectivity should not be a prerequisite for testing a new constitutive
law.  A staged study can begin with the present local eligibility test and isolate the
rate-law effect.

\section{Discrete algorithms worth testing}

\subsection{Algorithm A: lagged rate source}

This is the lowest-cost experiment.
\begin{enumerate}
  \item Advect VOF to obtain \(\alpha^*\) and compute the current eligibility mask.
  \item Evaluate a rate \(R(\alpha^*,p^n,\uv^*)\).
  \item Impose \(\divg\uv^{n+1}=-R/(1-\alpha^*)\) as a known source in the existing
        projection.
  \item Limit the predicted removal to available void and, if needed,
        \(\delta\alpha_c\le r_{\max}\alpha^*\).
  \item Update VOF using the same corrected fluxes used by continuity.
\end{enumerate}
Algorithm A can compare fixed-time, saturated-pressure, and Rayleigh-inspired rates with
minimal pressure-solver changes.  Lagged pressure may delay the onset by one step and can
be unstable for a strong pressure response.

\subsection{Algorithm B: implicit linear compliance}

Insert Eq.~\eqref{eq:simple-compliance} directly into the projection to obtain
Eq.~\eqref{eq:reaction-poisson}, or its incremental Truchas form
Eq.~\eqref{eq:reaction-poisson-incremental}.  Lag \(\alpha\), but solve updated pressure
and collapse simultaneously.  This is the cleanest first implementation of the
disappearing-boundary idea.  It should be compared against Algorithm A at weak coupling
and tested for conditioning as \(\alpha\to1\).

The capillary-biased signed variant uses Eq.~\eqref{eq:capillary-incremental}.  It has
the same matrix as the unbiased signed scheme and changes only the known pressure offset
on the right-hand side.  This makes it the most direct extension of the tested linear
prototype.

After the pressure solve, compute the collapse volume from either the integrated corrected
face fluxes or the cell divergence.  These two evaluations should agree to the pressure
solver tolerance.  That comparison is an important implementation diagnostic.

\subsection{Algorithm C: one-sided or nonlinear implicit response}

Use either the active-set form
\begin{equation}
  \divg\uv=-C(\alpha)\pos{p-p_v}
\end{equation}
or an inertial response derived from Eq.~\eqref{eq:inertial-rate}.  The pressure operator
is then nonlinear.  This algorithm is justified only after the linear-compliance version
has demonstrated that pressure coupling cures the stationary-pocket failure.

\section{Conservation, stability, and interpretation}

\subsection{Liquid-volume conservation}

The purpose of Eq.~\eqref{eq:general-divergence} is to make the increase in liquid volume
equal the net liquid inflow.  In a finite-volume implementation, the same face fluxes must
be used in the pressure correction and the VOF update.  A post hoc decrement of
\(\alpha\) would manufacture liquid and should be avoided.

Truchas documentation reports that larger relaxation values in the experimental model
can increase mass loss and that additional VOF subcycling improves the result
\cite{truchas-manual}.  This makes discrete flux consistency and subcycle coupling first-
class acceptance criteria for any successor.

\subsection{What is and is not physical}

The pressure-compliance model is best described as a numerical closure for a missing
subgrid pressure boundary.  Unless its parameters are derived and validated for a
specific gas-removal mechanism, it should not be presented as a physical bubble model.
Rayleigh scaling provides a defensible inertial time scale, but a cell-scale nonspherical
fragment adjacent to walls is far outside Rayleigh's assumptions.

The model is intentionally dissipative in the sense that it removes an unresolved degree
of freedom.  A useful study should monitor pressure work associated with
\(p\,\divg\uv\), because overly rapid annihilation may create spurious pressure transients
or remove kinetic energy at a grid-dependent rate.

\subsection{Pressure sign, interface motion, and cavitation}

A symmetric linear compliance permits positive divergence when \(p<p_v\).  In an
eligible mixed cell this can represent expansion of an existing unresolved interface,
especially when the cell is mostly void; it is not automatically nucleation.  Because
\(C(0)=0\), the closure cannot create void from an exactly full liquid cell.  It can,
however, amplify a tiny numerical remnant under tension, so thresholds, topology, and
pressure histories still require scrutiny.

The strict positive-part law forbids expansion but may make a mostly void cell behave
less like the open void that it is replacing.  Equation~\eqref{eq:capillary-signed}
offers an intermediate interpretation: reversible response near the large-void limit,
with an increasingly negative expansion threshold as the remnant shrinks.  It is not a
complete cavitation model because it contains neither a nucleation criterion nor gas
mass, thermodynamics, or bubble dynamics.

\section{Verification and discriminating tests}

The following hierarchy is designed to distinguish constitutive errors from topology and
VOF-transport errors.

\subsection{Local and one-dimensional tests}
\begin{enumerate}
  \item \textbf{Prescribed-pressure decay.}  Hold \(p-p_v\) fixed and verify the analytic
        \(\alpha(t)\) for fixed-time and linear-pressure relaxation over a sequence of
        time steps.
  \item \textbf{Finite-extinction law.}  Verify that an \(\alpha^{2/3}\) rate gives a
        linear decrease of \(\alpha^{1/3}\) in time and the predicted extinction time.
  \item \textbf{Reaction-Poisson cell.}  In a one-dimensional column, verify the smooth
        transition from a strong \(p_v\) anchor to an ordinary interior pressure as
        \(\alpha\) decreases.
  \item \textbf{Flux identity.}  Check that integrated net inflow equals the computed
        loss of void to solver tolerance.
  \item \textbf{Signed capillary threshold.}  Prescribe pressure on both sides of
        \(p_v-\Pi_\sigma\) and verify the predicted sign of divergence, including zero
        response at the equilibrium threshold.
\end{enumerate}

\subsection{Two- and three-dimensional tests}
\begin{enumerate}
  \item \textbf{Pressure-fed ceiling fill.}  This is the primary regression.  Record the
        top-layer void volume, pressure, divergence, and inflow after macroscopic filling
        stops.  The current model is expected to plateau; a successful pressure model
        should continue toward zero with pressure rising smoothly toward neighboring
        liquid pressure and divergence decreasing toward zero.
  \item \textbf{Stationary enclosed fragment.}  Initialize a trapped fragment with
        \(\Delta\alpha=0\) and positive surrounding pressure.  This directly tests
        persistent activation.
  \item \textbf{Translating fragment.}  Advect an enclosed fragment at nearly uniform
        velocity with weak pressure loading.  This exposes the spurious translation
        sensitivity of \(h/|\uv|\) scaling.
  \item \textbf{Exterior free surface and thin filament.}  Confirm that eligibility
        logic never collapses exterior-connected void, including underresolved
        connections near the classification threshold.
  \item \textbf{Multi-cell component.}  Compare local and graph-based eligibility and,
        for inertial laws, cell-size versus component-size length scales.
  \item \textbf{Capillary-bias sweep.}  Vary \(\alpha\), \(h\), \(\sigma\), and
        \(c_\sigma\).  Verify that large-void behavior approaches unbiased signed
        compliance, small remnants resist expansion, and the measured threshold scales
        as \(c_\sigma\sigma/h\).
\end{enumerate}

\subsection{Convergence matrix and diagnostics}

Each discriminating case should vary \(\Dt\), VOF subcycling, mesh spacing, collapse
parameters, and pressure-solver tolerance.  At minimum record
\begin{equation*}
  V_v(t)=\sum_i\alpha_iV_i,\quad
  \sum_i V_i\divg\uv_i,\quad
  p_{\min},p_{\max},\quad
  \text{liquid-volume error},\quad
  \text{pressure iterations}.
\end{equation*}
For the ceiling-fill problem, cellwise histories of \((\alpha,p,\divg\uv,C)\) in the top
layer are more informative than images alone.

\section{Recommended research sequence}

The models should be introduced in an order that makes failures interpretable.

\begin{enumerate}
  \item \textbf{Document the baseline.}  Record the exact Truchas commit, candidate-cell
        test, source formula, sign convention, and mapping between projection divergence
        and VOF correction.  Reproduce the ceiling-fill plateau quantitatively.
  \item \textbf{Implement a fixed-time source.}  Use Eq.~\eqref{eq:fixed-time} with the
        current eligibility test.  This establishes that persistent activation alone can
        eliminate the plateau and provides a time-step-convergence baseline.
  \item \textbf{Implement lagged saturated pressure relaxation.}  Use
        Eq.~\eqref{eq:pressure-relaxation}.  Determine whether pressure response improves
        the temporal and spatial pattern of collapse.
  \item \textbf{Characterize the signed-compliance prototype.}  The reaction term in
        Eq.~\eqref{eq:reaction-poisson-incremental} has been implemented and has performed
        well in limited testing.  Quantify pressure smoothness, conservation,
        conditioning, time-step dependence, and sensitivity to tiny residual voids.
  \item \textbf{Add the capillary bias.}  Implement
        Eq.~\eqref{eq:capillary-incremental} first with
        \(B(\alpha)=(1-\alpha)^2\).  Compare it directly with unbiased signed compliance and the
        strict one-sided active-set form.
  \item \textbf{Improve eligibility only after the rate law is understood.}  Add graph
        connectivity and hysteresis, using cases that distinguish multi-cell trapped void
        from exterior-connected filaments.
  \item \textbf{Evaluate inertial scaling.}  Test Eq.~\eqref{eq:inertial-rate} with both
        cell and component length scales.  Retain it only if it reduces parameter and mesh
        sensitivity relative to the simpler pressure law.
\end{enumerate}

The most promising next model is therefore
\begin{equation}
  \boxed{
  \divg\uv=-
  \trapped\,
  \frac{\alpha}{(1-\alpha)\tau_0p_0}
  \left[p-p_v+c_\sigma\frac{\sigma}{h}B(\alpha)\right],}
  \label{eq:recommended}
\end{equation}
treated implicitly in pressure with \(\alpha\) lagged and with a controlled pure-void
handoff.  The strict positive-part and unbiased signed laws are the essential
comparisons.  Equation \eqref{eq:recommended} preserves the successful linear structure
of the prototype, directly fixes the stationary-pocket failure, and embodies the desired
transition from reversible open-void response to capillary-biased disappearance of a
small remnant.

\section{Open questions}

The study should not hide the following unresolved choices.
\begin{enumerate}
  \item What fixed form of \(B(\alpha)\) provides the desired transition without adding
        poorly identifiable shape parameters?
  \item Should \(\sigma\) be the physical material property with only \(c_\sigma\)
        calibrated, or should the model expose a single effective parameter
        \(\Sigma_{\mathrm{eff}}\)?
  \item At what void fraction and pressure should positive divergence be permitted, and
        what safeguards prevent amplification of numerical remnants?
  \item What pressure should be used for \(p_v\): zero gauge or an imposed exterior
        pressure?
  \item Can \(\tau_0p_0\) be related to a discrete inertial impedance, reducing the number
        of user parameters?
  \item Should the compliance be singular as \(F\to0\), or merely large relative to the
        local pressure stencil?
  \item Is a component-aware length scale necessary for mesh convergence, or is deliberate
        cell-scale annihilation the more honest numerical model?
  \item How often must connectivity be recomputed in a distributed calculation, and what
        hysteresis thresholds prevent chatter without retaining false connections?
  \item Does pressure-compliance work create unacceptable transients when a cell changes
        from exterior-connected to trapped?
  \item How should collapse couple to VOF subcycling so that liquid volume remains
        conservative at both the flow-step and subcycle levels?
\end{enumerate}

\section{Conclusion}

The trapped-void pathology is best understood as the loss of an internal void-pressure
boundary when a fragment becomes underresolved.  The existing Truchas model addresses
the problem in the right place---the pressure projection---and plausibly tries to match
the intrinsic local interface-evolution time.  Its \(|\Delta\alpha|\) sensor, however,
vanishes precisely in the important case of a stationary pressure-loaded pocket.

All candidate models can be expressed as localized negative divergence consistent with
real-fluid conservation.  Fixed-time relaxation is the essential baseline; flow-time
laws preserve part of the old motivation but fail in stagnation; pressure and Rayleigh
laws retain a nonzero physical driver; and a true gas model answers a different question.
The pressure-compliance formulation unifies the most useful features.  It replaces an
abrupt switch between \(p=p_v\) and ordinary interior pressure with a reaction term whose
strength decreases with void fraction.  This is both a concrete algorithm and a clear
interpretation: a continuously disappearing subgrid pressure boundary.

Limited tests of the linear signed formulation support that interpretation.  A
capillary-biased signed law is a natural refinement: it preserves reversible response
for a mostly void cell, retains the same linear incremental pressure matrix, and shifts
the expansion threshold downward as the fragment becomes small.  Its value will depend
on choosing and validating the subgrid curvature function without duplicating resolved
surface-tension physics.

\appendix
\section{Dimensions and notation}

\begin{table}[htbp]
\centering
\small
\begin{tabularx}{\textwidth}{@{}p{0.15\textwidth}p{0.26\textwidth}X@{}}
\toprule
Symbol & Dimensions & Meaning \\
\midrule
\(F\) & 1 & real-fluid volume fraction \\
\(\alpha\) & 1 & void volume fraction \\
\(R\) & \(T^{-1}\) & material void-destruction rate \\
\(C\) & \(P^{-1}T^{-1}\) & pressure compliance in Eq.~\eqref{eq:compliance} \\
\(\lambda\) & \(P^{-1}T^{-1}\) & pressure-to-void-rate coefficient \\
\(p_v\) & \(P\) & nominal void pressure \\
\(p_{\mathrm{eq}}\) & \(P\) & capillary-shifted liquid-side equilibrium pressure \\
\(\Pi_\sigma\) & \(P\) & subgrid capillary pressure shift \\
\(\sigma\) & \(PL\) & physical liquid--void surface tension \\
\(c_\sigma\) & 1 & subgrid curvature/geometry coefficient \\
\(B(\alpha)\) & 1 & prescribed void-fraction activation function for capillary bias \\
\(p_0\) & \(P\) & pressure scale for relaxation \\
\(\tau_0\) & \(T\) & collapse relaxation time \\
\(h\) & \(L\) & representative cell length \\
\(V_c\) & \(L^3\) & cell volume in three dimensions \\
\bottomrule
\end{tabularx}
\end{table}

\section{Minimal pseudocode for incremental implicit compliance}

\begin{verbatim}
alpha_star = advect_vof(...)
trapped    = classify_trapped_void(alpha_star, ...)

for each flow cell i:
    if trapped(i) and F(i) > F_D:
        C(i) = alpha_star(i) / (F(i) * tau0 * p0)
        C(i) = min(C(i), C_max)
        h(i) = cell_volume(i) ** (1/3)       ! 3-D definition
        Pi_sigma(i) = c_sigma * sigma / h(i) * B(alpha_star(i))
    else:
        C(i) = 0
        Pi_sigma(i) = 0

! p_n is the existing pressure used in the velocity predictor
u_star = velocity_predictor_using(p_n)

! Capillary-biased signed compliance: one linear correction solve
assemble [-div(K grad) + C] phi
         = -div(u_star) - C*(p_n - p_v + Pi_sigma)
solve for phi

! For strict one-sided compliance instead, use an active-set iteration
! with active(i) = trapped(i) and p_n(i) + phi_iter(i) > p_v.

u_new = u_star - K * grad(phi)
p_new = p_n + phi
update VOF with pressure-corrected face fluxes
verify net liquid inflow = loss of void
\end{verbatim}

\begin{thebibliography}{10}

\bibitem{truchas-manual}
Computational Physics and Methods Group,
\emph{Truchas Reference Manual}, Version 21.01.1,
Los Alamos National Laboratory, 2021.
\url{https://www.truchas.org/ReferenceManual.pdf}.

\bibitem{truchas-source}
Truchas developers,
\emph{Truchas source repository}, Los Alamos National Laboratory.
\url{https://gitlab.com/truchas/truchas}.
Exact commit used for the source-level reconstruction remains to be recorded.

\bibitem{hirt-nichols1981}
C.~W. Hirt and B.~D. Nichols,
``Volume of Fluid (VOF) Method for the Dynamics of Free Boundaries,''
\emph{Journal of Computational Physics}, 39(1), 201--225, 1981.
\doi{10.1016/0021-9991(81)90145-5}.

\bibitem{bell-colella-glaz1989}
J.~B. Bell, P.~Colella, and H.~M. Glaz,
``A Second-Order Projection Method for the Incompressible Navier--Stokes Equations,''
\emph{Journal of Computational Physics}, 85(2), 257--283, 1989.
\doi{10.1016/0021-9991(89)90151-4}.

\bibitem{rayleigh1917}
Lord Rayleigh,
``On the Pressure Developed in a Liquid during the Collapse of a Spherical Cavity,''
\emph{The London, Edinburgh, and Dublin Philosophical Magazine and Journal of Science},
34(200), 94--98, 1917.
\doi{10.1080/14786440808635681}.

\bibitem{brennen1995}
C.~E. Brennen,
\emph{Cavitation and Bubble Dynamics}, Oxford University Press, 1995.

\bibitem{li-carrica2018}
J.~Li and P.~M. Carrica,
``An Approach to Couple Velocity/Pressure/Void Fraction in Two-Phase Flows with
Incompressible Liquid and Compressible Bubbles,''
\emph{International Journal of Multiphase Flow}, 102, 77--94, 2018.
\doi{10.1016/j.ijmultiphaseflow.2018.01.021}.

\bibitem{shin2020}
S.~Shin,
``Simulation of Compressibility of Entrapped Air in an Incompressible Free Surface Flow
Using a Pressure-Based Method for Unified Equations,''
\emph{International Journal for Numerical Methods in Fluids}, 92(10), 1274--1289, 2020.
\doi{10.1002/fld.4827}.

\bibitem{brackbill1992}
J.~U. Brackbill, D.~B. Kothe, and C.~Zemach,
``A Continuum Method for Modeling Surface Tension,''
\emph{Journal of Computational Physics}, 100(2), 335--354, 1992.
\doi{10.1016/0021-9991(92)90240-Y}.

\end{thebibliography}

\end{document}
